% !TEX root = ../Report.tex
\section{Methodology}
\label{sec:methodology}
\begin{figure*}[t]
\centering
\includegraphics[width=0.82\textwidth]{images/chapter5_overall_architecture.pdf}
\caption{KGE-NFM framework with alternative DistMult and CompGCN encoders. The KG encoder produces drug and protein embeddings that are combined with molecular and sequence descriptors before NFM prediction.}
\Description{A pipeline from biomedical knowledge graph to DistMult or CompGCN encoder, drug and protein embeddings, descriptor features, neural factorization machine, and DTI probability.}
\label{fig:kge_nfm_architecture}
\end{figure*}

\subsection{Framework Overview}

The proposed methodology follows the KGE-NFM paradigm for DTI prediction, in which relational biomedical evidence is first encoded at the entity level and then integrated with intrinsic drug and protein descriptors for supervised pair classification. The central design choice in this study is the knowledge graph encoder. The original KGE-NFM formulation uses DistMult to obtain drug and protein embeddings, whereas this work evaluates whether replacing that bilinear encoder with CompGCN improves the representations supplied to the downstream predictor. All other stages, including fold construction, descriptor processing, feature-field design, and the NFM classifier, are held fixed so that the comparison focuses on the effect of the KG representation model.

As shown in Figure~\ref{fig:kge_nfm_architecture}, the framework is trained sequentially rather than end-to-end. For each fold, the KG encoder is trained on supporting biomedical triples together with the positive DTI triples from the training split; test DTI labels are excluded from graph construction. The trained encoder then provides separate graph-derived embeddings for the drug and protein in each candidate pair. These embeddings are treated as complementary context features, not as replacements for entity descriptors. The final predictor therefore receives four aligned fields: drug KG embedding, protein KG embedding, drug Morgan fingerprint, and protein CTD descriptor. NFM projects these heterogeneous fields into a shared latent space and learns their pairwise and nonlinear interactions before producing the DTI probability.


\subsection{Problem Formulation and Training Graph}

For a drug $d$ and protein $p$, KGE-NFM estimates
\begin{equation}
    \hat{y}_{dp}=F(d,p),
\end{equation}
where $\hat{y}_{dp}\in[0,1]$ is the predicted interaction probability and the supervised label is $y_{dp}\in\{0,1\}$. The KGE stage learns representations through an encoder $g_{\theta}$:
\begin{equation}
    \mathbf{z}_{d}=g_{\theta}(d,\mathcal{G}),
    \qquad
    \mathbf{z}_{p}=g_{\theta}(p,\mathcal{G}),
\end{equation}
where $g_{\theta}$ is instantiated as either DistMult or CompGCN. The NFM stage then predicts
\begin{equation}
    \hat{y}_{dp}
    =
    \operatorname{NFM}
    \left(
        \mathbf{z}_{d},
        \mathbf{z}_{p},
        \mathbf{s}_{d},
        \mathbf{s}_{p}
    \right),
\end{equation}
with drug descriptor $\mathbf{s}_{d}$ and protein descriptor $\mathbf{s}_{p}$.

For each cross-validation fold, the KGE training graph is constructed as
\begin{equation}
    \mathcal{T}_{train}^{KGE}
    =
    \mathcal{T}_{support}
    \cup
    \mathcal{T}_{DTI,train}^{+}.
    \label{eq:kge_training_graph}
\end{equation}
Here, $\mathcal{T}_{support}$ denotes biomedical context triples such as pathway, annotation, domain, family, and interaction relations, while $\mathcal{T}_{DTI,train}^{+}$ contains only positive DTI triples from the training split. Negative triples are generated by corrupting heads or tails, and both KGE encoders use the margin-ranking objective in Equation~\ref{eq:background_kge_loss}. Feature scaling and optional PCA are fitted on training data and then applied to test data to prevent leakage.

\subsection{DistMult Encoder}

DistMult assigns a vector to each entity and relation \cite{yang2015embedding}. For a triple $(h,r,t)$,
\begin{equation}
    \mathbf{h},\mathbf{r},\mathbf{t}
    \in
    \mathbb{R}^{k}.
\end{equation}
The triple score is a diagonal bilinear interaction:
\begin{equation}
    f_{\mathrm{DM}}(h,r,t)
    =
    \mathbf{h}^{\top}
    \operatorname{diag}(\mathbf{r})
    \mathbf{t}
    =
    \sum_{i=1}^{k}h_i r_i t_i.
    \label{eq:distmult_score}
\end{equation}
This form is efficient because it avoids a full relation-specific matrix. After KGE training, drug and protein representations are read from the learned entity embedding table:
\begin{equation}
    \mathbf{z}_{d}^{DM}=\mathbf{E}[d],
    \qquad
    \mathbf{z}_{p}^{DM}=\mathbf{E}[p].
\end{equation}

DistMult is a useful baseline, but its score is symmetric for a fixed relation:
\begin{equation}
    f_{\mathrm{DM}}(h,r,t)
    =
    f_{\mathrm{DM}}(t,r,h).
\end{equation}
It also learns through triple-level scoring rather than explicit neighborhood aggregation. These limitations motivate the CompGCN encoder, which updates entity states using relation-conditioned messages.

\subsection{CompGCN Encoder}

CompGCN is a multi-relational graph convolutional encoder that jointly updates entity and relation representations \cite{vashishth2020composition}. For each original triple $(u,r,v)$, the graph is augmented with an inverse triple $(v,r^{-1},u)$ and self-loop relations. Original, inverse, and self-loop edges allow information to propagate with direction-aware semantics. At layer $k$, CompGCN maintains entity states $\mathbf{h}_{v}^{k}$ and relation states $\mathbf{h}_{r}^{k}$.

Before aggregation, a neighboring entity representation is composed with the relation embedding:
\begin{equation}
    \mathbf{m}_{u,r}^{k}
    =
    \phi
    \left(
        \mathbf{h}_{u}^{k},
        \mathbf{h}_{r}^{k}
    \right).
\end{equation}
The supported composition functions are subtraction, element-wise multiplication, and circular correlation:
\begin{align}
    \phi_{sub}(\mathbf{h}_{u},\mathbf{h}_{r})
    &=
    \mathbf{h}_{u}-\mathbf{h}_{r}, \\
    \phi_{mult}(\mathbf{h}_{u},\mathbf{h}_{r})
    &=
    \mathbf{h}_{u}\odot\mathbf{h}_{r}, \\
    \phi_{corr}(\mathbf{h}_{u},\mathbf{h}_{r})
    &=
    \mathbf{h}_{u}\star\mathbf{h}_{r}.
\end{align}
Here, $\odot$ denotes element-wise multiplication and $\star$ denotes circular correlation. The main configuration uses multiplication, so messages are relation-gated as
\begin{equation}
    \phi
    \left(
        \mathbf{h}_{u}^{k},
        \mathbf{h}_{r}^{k}
    \right)
    =
    \mathbf{h}_{u}^{k}\odot\mathbf{h}_{r}^{k}.
\end{equation}

The node update aggregates composed messages:
\begin{equation}
    \mathbf{h}_{v}^{k+1}
    =
    \sigma
    \left(
        \sum_{(u,r)\in\mathcal{N}(v)}
        \mathbf{W}_{\lambda(r)}^{k}
        \phi
        \left(
            \mathbf{h}_{u}^{k},
            \mathbf{h}_{r}^{k}
        \right)
    \right),
    \label{eq:compgcn_node_update}
\end{equation}
where $\mathcal{N}(v)$ is the relational neighborhood of $v$, $\sigma$ is a nonlinear activation, and $\lambda(r)\in\{O,I,S\}$ selects transformations for original, inverse, and self-loop edges. Thus, CompGCN shares direction-specific matrices $\mathbf{W}_{O}^{k}$, $\mathbf{W}_{I}^{k}$, and $\mathbf{W}_{S}^{k}$ rather than learning a separate full transformation for every relation. Relation embeddings are updated as
\begin{equation}
    \mathbf{h}_{r}^{k+1}
    =
    \mathbf{W}_{rel}^{k}
    \mathbf{h}_{r}^{k}.
    \label{eq:compgcn_relation_update}
\end{equation}
Stacking $L$ layers yields entity states $\mathbf{h}_{v}^{L}$ that encode relation-conditioned neighborhood information. Figure~\ref{fig:compgcn_message_passing} illustrates this message-passing mechanism.

\begin{figure}[t]
\centering
\includegraphics[width=\columnwidth]{images/chapter5_compgcn_message_passing.pdf}
\caption{CompGCN message passing. Neighboring entity states are composed with relation states and transformed according to edge direction before aggregation.}
\Description{A CompGCN message-passing diagram showing original, inverse, and self-loop relation-aware messages flowing into an entity update.}
\label{fig:compgcn_message_passing}
\end{figure}

CompGCN acts as an encoder whose enriched entity and relation states are passed to a link-prediction decoder:
\begin{equation}
    f_{CG}(h,r,t)
    =
    \operatorname{Dec}
    \left(
        \mathbf{h}_{h}^{L},
        \mathbf{h}_{r}^{L},
        \mathbf{h}_{t}^{L}
    \right).
\end{equation}
The encoder--decoder model is optimized with the same negative-sampling and margin-ranking setup used for DistMult. After training, the enriched drug and protein embeddings are extracted as
\begin{equation}
    \mathbf{z}_{d}^{CG}=\mathbf{h}_{d}^{L},
    \qquad
    \mathbf{z}_{p}^{CG}=\mathbf{h}_{p}^{L}.
\end{equation}
These vectors replace the DistMult embeddings in the downstream NFM while leaving the rest of the pipeline unchanged.

\subsection{Feature Construction}

For either encoder, the candidate pair contains four feature fields:
\begin{equation}
    \mathcal{F}_{dp}
    =
    \left\{
        \mathbf{z}_{d},
        \mathbf{z}_{p},
        \mathbf{s}_{d},
        \mathbf{s}_{p}
    \right\}.
    \label{eq:kge_nfm_fields}
\end{equation}
The graph fields $\mathbf{z}_{d}$ and $\mathbf{z}_{p}$ represent relational biomedical context. The descriptor fields are a 1024-dimensional Morgan fingerprint $\mathbf{s}_{d}\in\mathbb{R}^{1024}$ and a 147-dimensional protein CTD descriptor $\mathbf{s}_{p}\in\mathbb{R}^{147}$. Conceptually, the pair representation is
\begin{equation}
    \mathbf{x}_{dp}
    =
    [
        \mathbf{z}_{d}
        \mathbin{\Vert}
        \mathbf{z}_{p}
        \mathbin{\Vert}
        \mathbf{s}_{d}
        \mathbin{\Vert}
        \mathbf{s}_{p}
    ],
\end{equation}
but the implemented NFM processes the four blocks as separate fields. This preserves the distinction between drug graph context, protein graph context, molecular structure, and protein sequence descriptors.

\subsection{Neural Factorization Machine}

The downstream predictor is a field-based NFM derived from the NFM principle in \cite{he2017neural}. Each input field $\mathbf{x}_{dp}^{(m)}\in\mathbb{R}^{D_m}$ is projected into a shared latent dimension:
\begin{equation}
    \mathbf{u}_{m}
    =
    \mathbf{P}_{m}\mathbf{x}_{dp}^{(m)}+\mathbf{b}_{m}.
    \label{eq:nfm_field_projection}
\end{equation}
The projected fields are combined using Bi-Interaction pooling:
\begin{equation}
    \mathbf{f}_{BI}
    =
    \frac{1}{2}
    \left[
        \left(
            \sum_{m=1}^{M}\mathbf{u}_{m}
        \right)^2
        -
        \sum_{m=1}^{M}\mathbf{u}_{m}^{2}
    \right].
    \label{eq:nfm_bi_interaction}
\end{equation}
This operation summarizes pairwise interactions among the four fields. The pooled vector is then processed by hidden layers:
\begin{equation}
    \mathbf{a}^{(l+1)}
    =
    \operatorname{ReLU}
    \left(
        \mathbf{W}^{(l)}\mathbf{a}^{(l)}+\mathbf{b}^{(l)}
    \right),
    \qquad
    \mathbf{a}^{(0)}=\mathbf{f}_{BI}.
\end{equation}

The final logit combines field-level linear terms and the nonlinear interaction network:
\begin{equation}
    o_{dp}
    =
    \sum_{m=1}^{M}
    \left(
        \mathbf{a}_{m}^{\top}\mathbf{x}_{dp}^{(m)}+c_m
    \right)
    +
    \operatorname{MLP}(\mathbf{f}_{BI}).
\end{equation}
The predicted DTI probability is
\begin{equation}
    \hat{y}_{dp}=\operatorname{sigmoid}(o_{dp}).
\end{equation}
NFM is optimized with binary cross-entropy:
\begin{equation}
    \mathcal{L}_{DTI}
    =
    -\frac{1}{N}
    \sum_{i=1}^{N}
    \left[
        y_i\log(\hat{y}_i)
        +(1-y_i)\log(1-\hat{y}_i)
    \right].
\end{equation}
This final stage is where relational embeddings and descriptor fields are jointly converted into an interaction probability.
